% Fragmento integrable. Preambulo anfitrion: amsmath, amssymb, longtable, hyperref.
\section{Transporte ordenado y representación de Jacobi}
\label{hmtfg:fragmento:transporte-ordenado-y-representacion-de-jacobi}

\subsection{Transporte completo desde el continuo y realización espectral del perfil}
\label{hmtfg:escalado:14}

El transporte ordenado del continuo y la representación espectral del lector se componen antes de la selección paramétrica.

\subsubsection{Cobertura de todas las raíces y transporte de su mínimo}
\label{hmtfg:escalado:14.1}

La construcción ordenada del continuo desarrollada en los antecedentes produce la publicación arquimediana sobreyectiva desde cilindros compatibles y su cociente ordenado. Dentro de su universo declarado \(\mathfrak G\), demuestra el isomorfismo de cuerpos ordenados completos
\begin{equation}
\iota:\mathbb R_{\rm HMT}^{\mathfrak G}
   \xrightarrow{\cong}\mathbb R^{\mathfrak G}.
\label{hmtfg:escalado:eq:37}
\end{equation}
Se cambia aquí el nombre \(\eta\) del propietario por \(\iota\), para distinguirlo del parámetro de ponderación, fijado a cero. La segunda proyección conserva incidencia, frontera y memoria sobre el mismo antecedente; el cociente de evaluación y la historia completa mantienen sus dominios respectivos.

Construyamos en el cuerpo HMT, antes de seleccionar una raíz,
\[
\cos_{\rm H}(t)=\sum_{k=0}^{\infty}\frac{(-1)^kt^{2k}}{(2k)!},
\qquad
u_{\rm H}(x)=\frac1{12}\sum_{m=1}^{12}
\left[1-\cos_{\rm H}\!\left(\frac{2(3m-1)x}{\sqrt[\rm H]{899}}\right)\right],
\]
\begin{equation}
F_a(x)=1-a\,u_{\rm H}(x),\qquad S_n^{\rm H}(a)=F_a^{\,2^n}(0).
\label{hmtfg:escalado:eq:38}
\end{equation}
Los enteros y el segundo momento son los ya derivados del lector dodecafásico. La raíz cuadrada es la positiva del cuerpo ordenado. La serie convergente realiza analíticamente ese lector; conserva el perfil y el coeficiente normalizador de sección~\ref{hmtfg:escalado:1}.

\textbf{Teorema de transporte del retorno.} Para todo parámetro del dominio interno de \eqref{hmtfg:escalado:eq:38},
\begin{equation}
\iota(F_a(x))=f_{\iota(a)}(\iota(x)),\qquad
\iota(S_n^{\rm H}(a))=S_n(\iota(a)).
\label{hmtfg:escalado:eq:39}
\end{equation}
\textbf{Demostración.} El isomorfismo ordenado fija los racionales y preserva la raíz positiva. Preserva los límites convergentes: cada entorno de radio racional positivo se transporta a otro del mismo radio. Aplicado a las sumas parciales de \eqref{hmtfg:escalado:eq:38}, transporta el coseno y \(u_{\rm H}\). La primera igualdad resulta de conservar suma y producto; la segunda, por inducción sobre el número de iteraciones.

Fijado un centro anterior \(a_{n-1}\), definamos \(Z_n^{\rm H}\) por retorno cero, período crítico exacto \(2^n\) y \(a>a_{n-1}\), dentro del dominio paramétrico declarado. Sea \(Z_n\) el conjunto definido por las mismas condiciones para \(f_\lambda\), a la derecha de \(\iota(a_{n-1})\). Entonces
\begin{equation}
\boxed{\iota[Z_n^{\rm H}]=Z_n,\qquad
\iota(\min Z_n^{\rm H})=\min Z_n.}
\label{hmtfg:escalado:eq:40}
\end{equation}
La segunda igualdad se afirma cuando existe el mínimo; la primera preserva también la no existencia de candidatos.

\textbf{Demostración.} \eqref{hmtfg:escalado:eq:39} preserva el retorno cero y los retornos anteriores que determinan el período exacto. El orden conserva la posición respecto del centro previo. La sobreyectividad de \eqref{hmtfg:escalado:eq:37}, aplicada a cualquier parámetro raíz, prueba la cobertura inversa. Si \(a_*\) es el mínimo, cualquier \(z\in Z_n\) tiene la forma \(\iota(a)\), con \(a\in Z_n^{\rm H}\), de donde \(\iota(a_*)\le z\). La afirmación inversa utiliza \(\iota^{-1}\).

La sobreyectividad cubre todos los parámetros de la carta, incluidos aquellos cuyas raíces no se hayan aislado por un cálculo finito. El isomorfismo conserva el universo de interpretación declarado y el orden de selección.

\subsubsection{Representación de Jacobi del perfil dodecafásico completo}
\label{hmtfg:escalado:14.2}

Definamos la medida atómica del lector,
\begin{equation}
s_m=\frac{4(3m-1)^2}{899},\qquad
\nu_0=\frac1{12}\sum_{m=1}^{12}\delta_{s_m},\qquad
M_r=\int s^r\,d\nu_0(s).
\label{hmtfg:escalado:eq:41}
\end{equation}
Los doce nodos son distintos y positivos. Para un polinomio \(p\ne0\) de grado menor que \(r\le12\),
\begin{equation}
\sum_{i,j=0}^{r-1}p_iM_{i+j}p_j
=\frac1{12}\sum_{m=1}^{12}p(s_m)^2>0.
\label{hmtfg:escalado:eq:42}
\end{equation}
Un polinomio de grado menor que doce no se anula en los doce nodos. En grado doce aparece \(\prod_m(s-s_m)\), de norma nula.

La ortogonalización respecto de esta medida produce polinomios ortonormales \(\psi_0,\ldots,\psi_{11}\), con \(\psi_0=1\). Sean
\[
Q_{mj}=\psi_j(s_m)/\sqrt{12},\quad
D=\operatorname{diag}(s_m),\quad J_{12}=Q^{\mathsf T}DQ.
\]
Se cumplen
\begin{equation}
Q^{\mathsf T}Q=I,\quad DQ=QJ_{12},\qquad
\boxed{u_0(x)=e_0^{\mathsf T}
\bigl[I-\cos(x\sqrt{J_{12}})\bigr]e_0.}
\label{hmtfg:escalado:eq:43}
\end{equation}
\textbf{Demostración.} La ortogonalidad procede de \eqref{hmtfg:escalado:eq:42}. La multiplicación por \(s\) produce una recurrencia de tres términos: los elementos separados por más de una posición se anulan por ortogonalidad y grado. Fijar positivos los coeficientes principales hace positivas las subdiagonales. Por cálculo funcional,
\(\cos(x\sqrt{J_{12}})=Q^{\mathsf T}\cos(x\sqrt D)Q\).
El vector \(Qe_0\) es constante, de valor \(1/\sqrt{12}\); sustituirlo recupera la suma de sección~\ref{hmtfg:escalado:1}.

El artículo de Grassmannianos desarrolla este mecanismo para medidas marcadas y refinamientos. Aquí se aplica a \eqref{hmtfg:escalado:eq:41}; el rango doce pertenece a este lector. Las marcas genealógicas se conservan junto al lector, conforme a la fidelidad marcada del propietario. \(J_{12}\) representa el perfil analítico; el estado completo conserva adicionalmente rutas, orientación y memoria.

El verificador de esta ampliación comprueba con fracciones exactas los doce grados de norma positiva, la anulación del grado doce, el entrelazador \(VT=DV\) en base mónica y la autoadjunción ponderada \(HT=T^{\mathsf T}H\). Comprueba los momentos de grado cero a treinta; la identidad para todo grado se sigue del entrelazador. Sus primeros coeficientes son
\begin{equation}
a_0=2,\qquad \beta_1=\frac{2493348}{808201}.
\label{hmtfg:escalado:eq:44}
\end{equation}

\subsubsection{Sensibilidad orientada del retorno}
\label{hmtfg:escalado:14.3}

Para un retorno exacto de período \(p=2^n\), sean
\[
c_j=f_\lambda^j(0),\quad q_j=\partial_\lambda c_j,\quad
D_j=\prod_{k=1}^j f_\lambda'(c_k),\quad D_0=1.
\]
Antes del retorno crítico, los factores de \(D_{p-1}\) son distintos de cero. Diferenciación y telescopaje dan
\begin{equation}
q_{j+1}=-u_0(c_j)+f_\lambda'(c_j)q_j,\qquad
\boxed{\frac{q_p}{D_{p-1}}
=-\sum_{j=1}^{p-1}\frac{u_0(c_j)}{D_j}.}
\label{hmtfg:escalado:eq:45}
\end{equation}
En la palabra canónica, \(\operatorname{sign}c_j=(-1)^{v_2(j)}\). El signo de \(f_\lambda'(c_j)\) es el opuesto; por tanto,
\begin{equation}
\operatorname{sign}D_j=(-1)^{j+\sum_{k=1}^jv_2(k)}
=(-1)^{s_2(j)},
\label{hmtfg:escalado:eq:46}
\end{equation}
usando \(\sum_{k=1}^jv_2(k)=j-s_2(j)\), donde \(s_2(j)\) cuenta los unos binarios. La sensibilidad normalizada es
\begin{equation}
\mathcal T_n=-\sum_{j=1}^{2^n-1}(-1)^{s_2(j)}t_j,\qquad
t_j=\frac{u_0(c_j)}{|D_j|}>0.
\label{hmtfg:escalado:eq:47}
\end{equation}
La positividad nodal de \eqref{hmtfg:escalado:eq:42} determina \(u_0\); \eqref{hmtfg:escalado:eq:47} conserva además los productos orbitales y sus orientaciones.

La representación adicional por momentos exigiría una medida positiva \(\rho\) en \([0,1]\) con \(t_j=\int z^{j-1}\,d\rho\). Bajo esa condición,
\begin{equation}
\mathcal T_n=\int_0^1
\frac{1-\prod_{r=0}^{n-1}(1-z^{2^r})}{z}\,d\rho(z)>0
\label{hmtfg:escalado:eq:48}
\end{equation}
para \(\rho\ne0\), extendiendo el integrando por su valor \(1\) en cero. La identidad se obtiene expandiendo el producto binario.

El control racional en la raíz de período cuatro previamente certificada obtiene
\[
t_1\approx0.40425224267600139730,\quad
t_2\approx0.04925249167432646403,\quad
t_3\approx0.15740870098294833737
\]
y el intervalo exterior riguroso
\begin{equation}
\begin{gathered}
0.296096033367379523967764444238832345360107<\mathcal T_2,\\
\mathcal T_2<0.296096033367379523967764444238832345360112.
\end{gathered}
\label{hmtfg:escalado:eq:49}
\end{equation}
También certifica \(t_3>t_2\). Una medida positiva sobre \([0,1]\) impondría \(t_3\le t_2\); ese levantamiento concreto de los pesos orbitales a momentos positivos queda descartado. El signo positivo de \eqref{hmtfg:escalado:eq:49} permanece certificado. Este control distingue el fallo de una prueba propuesta del fallo de la propiedad que se pretendía probar.
